% FRAGMENT A13 (drafter A: class model and main results)
% REPLACES: main.tex lines 1399–1660 ("\subsection{A.2 Global trading and the main comparison}" to "… Online Appendix A.5 has the domain checks. \(\square\)")
%   This covers Appendix A.2: the global order bound, the proof of Proposition 2,
%   Proposition A.3, and Proposition A.4.
% LABELS KEPT: pa-proofs, the-global-order-bound, eq:paper-a-4 to eq:paper-a-14,
%   proof-of-proposition-2, deterrence-at-a-fixed-experiment, result-proposition-a-3,
%   threshold-distinctions, result-proposition-a-4.
% LABELS REMOVED: none here (the references to eq:paper-a1 at old lines 1476, 1494, 1538
%   disappear with this fragment).  LABELS ADDED: none.
% TAGS: (A.4)–(A.14) keep their numbers and objects. (A.6) now carries \kappa(\pi);
%   (A.7) gives J_H and J_L at a general prior (equal at \pi=1/2, as used by Appendix A.7);
%   (A.8) carries the prior factor and reduces to the old display at \pi=1/2;
%   (A.10)–(A.14) are unchanged and are stated for the point-mass benchmark.
% NEW SYMBOLS: J_H, J_L, K (in (A.7)); F_0, F_1, F_2 as in Proposition 2.
% NEW CITATIONS: none.  NUM PLACEHOLDERS: none.
% SOURCES: Q1_N_main.md Claims 4(a), 6–10 (proof steps), 8 (J_H, J_L, no-trade condition at a
%   general prior); SYNTHESIS.md A.1 rows 4–5.

\subsection{A.2 Global trading and the main comparison}\label{pa-proofs}

\subsubsection{The global order bound}\label{the-global-order-bound}

Fix a candidate pricing and preparation schedule. For a correctly signed
order of magnitude \(s\in[0,1]\) define

\[
\begin{aligned}
F_H(s)&=\int f(x-s)A_H(x)\,dx,\qquad
F_L(s)=\int f(x+s)A_L(x)\,dx,\\
U_j(s)&=sF_j(s)-ks.
\end{aligned}
\tag{A.4}\label{eq:paper-a-4}
\]

The residual \(A_j\) is bounded and measurable but need not be
continuous, since it jumps where entry jumps. The deviation enters only
through the translated density. For \(s_1<s_2\),

\[
F_j(s_2)-F_j(s_1)
=\int_{s_1}^{s_2}\int \partial_u f(x\mp u)\,A_j(x)\,dx\,du,
\tag{A.5}\label{eq:paper-a-5}
\]

The Laplace density is absolutely continuous. Fubini's theorem
applies, because the absolute double integral is at most
\(\|A_j\|_\infty\,|s_2-s_1|\,\|f'\|_1\). So \(F_j\) is absolutely
continuous, and \(|f'|\le f/b\) gives \(|F_j'|\le F_j/b\) almost
everywhere. A unilateral deviation never differentiates an entry jump.
The derivative falls on the density alone. With the residual lower
bound in \eqref{eq:paper-11},

\[
U_j'(s)=F_j(s)+sF_j'(s)-k
\ge\left(1-\frac sb\right)F_j(s)-k
\ge\left(1-\frac1b\right)\rho\kappa(\pi)\Delta_T-k>0
\tag{A.6}\label{eq:paper-a-6}
\]

under the strong-incumbent inequality in \eqref{eq:paper-a3}, for every
\(s\in[0,1]\) and both classes. Integrating the positive lower bound
shows that the full correctly signed order strictly dominates every
smaller magnitude; this is global optimality, not a first-order
condition. The class-specific floors are \(\rho[1-M(\pi)]\Delta_T\) for
\(A_H\) and \(\rho m(\pi)\Delta_T\) for \(A_L\); \(\kappa(\pi)\) is the
smaller of the two. Online Appendix A.3 retains the \(W^{1,1}\) and
null-set details.

A sharper existence test applies to a fixed full-order candidate.
Substituting Bayes' rule with prior \(\pi\) into both residuals gives

\[
\begin{aligned}
J_H&=F_H(1)=(1-\pi)K,\qquad J_L=F_L(1)=\pi K,\\
K&=\Delta_T\int e(x)\frac{f(x-1)f(x+1)}{\pi f(x-1)+(1-\pi)f(x+1)}\,dx,
\end{aligned}
\tag{A.7}\label{eq:paper-a-7}
\]

so at \(\pi=1/2\) the two coincide, \(J=F_H(1)=F_L(1)=K/2\). The
pointwise ratio \(f(x-s)/f(x-1)\ge e^{-(1-s)/b}\) gives
\(F_j(s)\ge e^{-(1-s)/b}J_j\). Since \((1-s/b)e^{-(1-s)/b}\) is
decreasing on \([0,1]\), the inequality \(k<(1-1/b)\min\{J_H,J_L\}\)
makes every correctly signed marginal profit positive against that
candidate. This verifies existence of the full-order profile; it says
nothing about other candidate schedules.

\subsubsection{\texorpdfstring{Proof of
\hyperref[result-proposition-2]{Proposition
2}}{Proof of Proposition 2}}\label{proof-of-proposition-2}

Steps 1 to 4 force the orders in every candidate equilibrium. Step 5
builds the price and preparation schedules at those orders.

\emph{Step 1: beliefs, the floor, and the payoff ranking.} Positivity of
\(f\) makes every conditional flow density positive, and
\hyperref[result-proposition-a-1]{Proposition A.1} gives
\(m(\pi)\le\mu_X\le M(\pi)\) under every trading strategy; the price
posterior obeys the same bound. The low-cost challenger prepares at
every price in every economy, so cost-averaged entry is at least
\(\rho\) everywhere. Under the restrictions of
\hyperref[sec-values-preparation-and-the-sale]{Section 2.2},
\eqref{eq:paper-6} gives \(g_H\ge g_L\) and \(\Delta_T\ge0\) in each
economy. Condition \eqref{eq:paper-a3} with \(k>0\) gives
\(\Delta_T(F_1)>0\).
\hyperref[result-proposition-1]{Proposition 1}(a) gives
\(B_{F_1}(\pi)\le B_{F_0}(\pi)\), so \eqref{eq:paper-a2} gives
\(B_{F_1}(\pi)<c_H<B_{F_1}(M(\pi))\); since \(M(\pi)>\pi\), this forces
\(g_H(F_1)>g_L(F_1)\). The strengthening from \(F_0\) to \(F_1\) is
used here.

\emph{Step 2: residuals.} In either economy entry depends on the price
and the independent cost, and the challenger's value depends on flow
only through its class. Competitive pricing therefore gives
\(\mathbb E[V_T\mid j,x]=t_0+e(P(x))(t_j-t_0)\) and the residuals
\(A_H=e(P(x))\Delta_T(1-\mu_X)\) and \(A_L=e(P(x))\Delta_T\mu_X\), with
\(\rho\le e(P(x))\le1\), as in the residual step of
\hyperref[result-proposition-a-2]{Proposition A.2}. In the strong
economy, where \(\Delta_T(F_1)>0\), Proposition A.2 also gives
\(\Pr(H\mid P)=\mu_X\) almost surely and entry \eqref{eq:paper-9}. This
uses the independence of \(R\) and \(C\) from trading, not an assumption
that the buyer sees order flow.

\emph{Step 3: the weak economy.} A correctly signed order of size
\(s>0\) earns \(s[F_j(s)-k]\le s[\Delta_T(F_0)-k]<0\) by the upper
bound \(A_j\le\Delta_T\) and the first inequality in
\eqref{eq:paper-a3}. A wrong-signed order has nonpositive gross payoff
and still pays its cost. Zero is therefore the only best response
against every candidate schedule, and a mixture placing positive
probability on nonzero orders averages strictly negative payoffs with
zero. Under zero orders \(X=Z\) and the posterior is \(\pi\);
\eqref{eq:paper-a2} gives \(B_{F_0}(\pi)<c_H\), so only the low-cost
challenger prepares and entry is \(\rho\). The constant price
\(t_0+\rho[\pi t_H+(1-\pi)t_L-t_0]\) is competitive. This constructs the
unique outcome of part (i).

\emph{Step 4: the strong economy.} Against every candidate schedule,
\eqref{eq:paper-a-6} holds on the whole order interval. Every
equilibrium therefore has \((q_H,q_L)=(1,-1)\), and mixing cannot
introduce another optimal action.

\emph{Step 5: existence.} Full orders give

\[
\mu_X(x)=\left[1+\frac{1-\pi}{\pi}
\exp\left\{-\frac{|x+1|-|x-1|}{b}\right\}\right]^{-1},
\tag{A.8}\label{eq:paper-a-8}
\]

equal to \(m(\pi)\) on \(x\le-1\), strictly increasing on \((-1,1)\),
and equal to \(M(\pi)\) on \(x\ge1\). Define entry and price by
\eqref{eq:paper-9}. The strict monotonicity \eqref{eq:paper-a-3} gives a
measurable inverse on the price's image, including across its upward
jump and at the two plateau atoms. The buyer recovers the posterior from
the price, prepares optimally, and bids truthfully;
\eqref{eq:paper-a-6} verifies investor optimality; the price is
competitive by construction.

\emph{Step 6: comparisons.} By Step 1, \(\tau\in(\pi,M(\pi))\). On
\((-1,1)\) the log odds of \eqref{eq:paper-a-8} are
\(\log[\pi/(1-\pi)]+2x/b\), so high-cost preparation occurs exactly when
\(X\ge x^*\), with \(x^*\) in \eqref{eq:paper-12} and
\(x^*\in(0,1)\). Given class \(H\), \(X=1+Z\); given class \(L\),
\(X=-1+Z\). The Laplace survival function evaluated at \(x^*-1\) and
\(x^*+1\) gives \(\alpha_H\) and \(\alpha_L\) in \eqref{eq:paper-12},
both positive, so entry exceeds \(\rho\). Given its class and
preparation, the challenger's value is independent of the price, and it
wins with probability \(W_j(F)=\mathbb E_j[F(\theta)\mathbf1\{\theta\ge
p\}]\). High-class ownership is therefore \(\pi\rho W_H(F_0)\) at
\(F_0\) and \(\pi e_HW_H(F_1)\) at \(F_1\), which gives the stated
equivalence. Because \(x^*\in(0,1)\), \(\alpha_H>1/2\) and
\(e_H>(1+\rho)/2\), so \((1+\rho)W_H(F_1)\ge2\rho W_H(F_0)>0\) implies
\(e_HW_H(F_1)>\rho W_H(F_0)\). In the point-mass benchmark \(h\)
exceeds both incumbents' supports, so \(W_H=1\) in both economies and
\(\mathsf O_H=e_H/2>\rho/2\). Applying a constant kernel to the
strong-economy price reproduces the weak economy's constant experiment,
while no state-independent kernel applied to a constant experiment can
produce the nonconstant state-dependent price law of the strong economy;
the strong experiment therefore strictly Blackwell dominates the weak
one. All inequalities are strict and continuous in the primitives, so
for uniform incumbents the region is open in \((r_0,r_1,c_H,k)\);
\hyperref[pa-calculations]{Appendix A.4} shows that it is nonempty for
every pair of classes with
\(\mathbb E_H[\max\{\theta,p\}]>\mathbb E_L[\max\{\theta,p\}]\).

\emph{Part (iii).} Under the restrictions of
\hyperref[sec-values-preparation-and-the-sale]{Section 2.2},
\hyperref[result-proposition-1]{Proposition 1}(b) gives
\(\Delta_T(F_2)\ge\Delta_T(F_1)\), so the trading bound at \(F_2\)
follows from \eqref{eq:paper-a3}. The low-cost floor and the derivative
bound give unique full orders exactly as in Steps 4 and 5, while
\(B_{F_2}(M(\pi))<c_H\) excludes every expensive entrant because
\(M(\pi)\) is the largest attainable belief under any strategy and
\(g_H(F_2)\ge g_L(F_2)\). Entry returns to \(\rho\). The condition on
\(F_2\) is separate from \eqref{eq:paper-a2} and \eqref{eq:paper-a3}
and establishes a rise and a subsequent fall across three economies, not
a monotone path. \(\square\)

\Needspace{10\baselineskip}

\subsubsection{Deterrence at a fixed
experiment}\label{deterrence-at-a-fixed-experiment}

\protect\phantomsection\label{result-proposition-a-3}{\textbf{Proposition
A.3 (deterrence at a fixed information experiment; analytical).}}
\emph{Let a signal \(S\) and the class have a fixed joint distribution,
let the challenger's value depend on \(S\) only through its class, and
let the preparation cost \(C\) be independent of the signal, the class,
and the value, with a fixed law, as in the model. Hold this experiment
and cost law fixed as the incumbent strengthens from \(F_0\) to a
first-order stochastic strengthening \(F_1\), so the buyer's posterior
\(\Pr(H\mid S)\) does not depend on the incumbent. Then entry is weakly
decreasing in incumbent strength. It decreases strictly when the decline
in gross profit crosses preparation costs on a positive-probability set
of signal and cost realizations.}

Write \(\mu(S)=\Pr(H\mid S)\); cost independence makes this the buyer's
complete posterior about the class, and its expected gross profit given
\(S\) is \(B(\mu(S))\) because its value depends on \(S\) only through
the class. By \hyperref[result-proposition-1]{Proposition 1}(a),
\(B_{F_1}(\mu)\le B_{F_0}(\mu)\) at every \(\mu\), for every pair of
classes, so

\[
\mathbf1\{C\le B_{F_0}(\mu(S))\}-\mathbf1\{C\le B_{F_1}(\mu(S))\}
=\mathbf1\{B_{F_1}(\mu(S))<C\le B_{F_0}(\mu(S))\}\ge0.
\tag{A.9}\label{eq:paper-a-9}
\]

Taking expectations proves the weak decrease and the exact strictness
condition, without an atomless cost distribution. The result applies
within a region of fixed full orders, where the conditional flow laws do
not depend on the incumbent; it does not apply across order profiles
that change with the incumbent, because changing orders change the
experiment. A version with cost correlated with the signal would
condition on \(\Pr(H\mid S,C)\) and is not needed here. \(\square\)

\Needspace{10\baselineskip}

\subsubsection{Threshold distinctions}\label{threshold-distinctions}

\protect\phantomsection\label{result-proposition-a-4}{\textbf{Proposition
A.4 (pooling, full orders, and expensive-entry feasibility;
analytical).}} \emph{In the point-mass benchmark, restrict attention to
the maintained domain
\(\mathcal D=\{r\in(\ell,h):B_r(1/2)<c_H\}\). Define}

\[
\begin{aligned}
\mathfrak r(d)&=\ell+d+\sqrt{d^2+2\ell d},\\
r_N&=\mathfrak r(2k/\rho),\qquad
r_U=\mathfrak r\left(\frac{k}{(1-1/b)\rho m}\right).
\end{aligned}
\tag{A.10}\label{eq:paper-a-10}
\]

\emph{No trade is an equilibrium exactly when
\(\rho\Delta_T(r)/2\le k\), that is, when \(r\le r_N\). The stronger
restriction \(r<\mathfrak r(k)\) guarantees that it is the unique
outcome. The restriction \(r>r_U\) guarantees unique full orders.
Expensive entry is impossible in every equilibrium when \(B_r(M)<c_H\).
If the equality \(B_r(M)=c_H\) has a solution in \(\mathcal D\), its
unique crossing is}

\[
r_C=\frac{Mh-c_H+\sqrt{(Mh-c_H)^2+M[(1-M)\ell^2-p^2]}}{M}.
\tag{A.11}\label{eq:paper-a-11}
\]

\emph{At \(r_C\) the tie rule of \hyperref[sec-model]{Section 2} is
retained.}

Under no trade the posterior is \(1/2\) and entry is \(\rho\), so a
unilateral correctly signed order of size \(s\) sees the constant
residual \(\rho\Delta_T/2\) and earns

\[
s\left(\frac{\rho\Delta_T(r)}{2}-k\right).
\tag{A.12}\label{eq:paper-a-12}
\]

Zero is optimal exactly when the coefficient is nonpositive. Solving
\(\Delta_T(r)=d\) gives \((r-\ell)^2=2rd\), whose root above \(\ell\) is
\(\mathfrak r(d)\); since \(\Delta_T\) is increasing, the
pooling-existence boundary is \(r_N\). Equality at \(r_N\) leaves the
investor indifferent among correctly signed sizes against that constant
schedule; it does not make every other profile an equilibrium with
recomputed prices. The sufficient uniqueness boundary \(\mathfrak r(k)\)
comes from \(\Delta_T<k\), which defeats any nonzero order against every
candidate schedule by Step 3 above. The full-order boundary \(r_U\)
comes from \eqref{eq:paper-a-6} at \(\pi=1/2\). The ceiling profit is

\[
B_r(M)=Mh-\frac{Mr}{2}+\frac{(1-M)\ell^2-p^2}{2r},
\tag{A.13}\label{eq:paper-a-13}
\]

strictly decreasing on \(r>\ell>p\), so it crosses \(c_H\) at most once;
\eqref{eq:paper-a-11} is the larger root of
\(Mr^2-2(Mh-c_H)r-[(1-M)\ell^2-p^2]=0\) and is the crossing when it lies
in \(\mathcal D\). Because \(M\) is the largest posterior under any
strategy, \(B_r(M)<c_H\) excludes expensive entry regardless of orders.
The four objects answer different questions: an existence boundary, two
sufficient uniqueness bounds, and a participation ceiling; none
identifies the first appearance of an informative equilibrium.

Equations \eqref{eq:paper-a-10}, \eqref{eq:paper-a-11}, and
\eqref{eq:paper-a-13} are point-mass formulas. The arguments extend.
With classes and prior \(\pi\), a unilateral order of size \(s\)
against the constant schedule at the prior earns
\(s[e(\pi)\Delta_T(1-\pi)-k]\) in class \(H\) and
\(s[e(\pi)\Delta_T\pi-k]\) in class \(L\), where \(e(\pi)\) is entry at
the prior. No trade is therefore an equilibrium exactly when
\(e(\pi)\Delta_T\max\{\pi,1-\pi\}\le k\), which is
\(\rho\Delta_T\max\{\pi,1-\pi\}\le k\) when \(B(\pi)<c_H\). The other
thresholds solve the same equations with \(\Delta_T\) and \(B\) from
\eqref{eq:paper-6}, \(M(\pi)\) in place of \(M\), and \(\kappa(\pi)\) in
place of \(m\).

Under full Laplace orders \(\mu_X=M\) on the positive-probability tail
\(x\ge1\), so indifference at \(r_C\) is not a null event and the tie
rule admits expensive entry on that tail. As \(r\uparrow r_C\), the
threshold \(x^*\to1\) and

\[
\alpha_H\to\frac12,\qquad\alpha_L\to\frac12e^{-2/b},\qquad
\mathsf E(r)\to\rho+\frac{1-\rho}{4}\left(1+e^{-2/b}\right),
\tag{A.14}\label{eq:paper-a-14}
\]

whereas strictly above \(r_C\) expensive entry is impossible. Under
logistic noise the posterior reaches \(M\) only in the limit, so
expensive entry converges to zero continuously
(\hyperref[pa-calculations]{Appendix A.4}). Online Appendix A.5 has the
domain checks. \(\square\)
